% =====================================================================
\chapter{Soluciones de los ejercicios}
\label{ap:soluciones}
% =====================================================================

Cada solución se presenta en cinco partes: datos, planteamiento, desarrollo, resultado y comprobación. En los pasos intermedios se redondea a cuatro cifras significativas salvo que se indique otra cosa. Se usan $e=2{,}718$, $e^{-1}=0{,}3679$, $e^{2}=7{,}389$ y $e^{-2}=0{,}1353$; 1~KiB $=1024$~bytes, 1~MiB $=1024^{2}$~bytes y 1~GiB $=1024^{3}$~bytes, mientras que GB y TB son potencias de~10.

% ---------------------------------------------------------------------
\section{Tema 1}

\begin{solucion}{eje:1-topk-topp}
\datos \emph{top-k} con $k=50$; núcleo con $p=0{,}9$. Caso (a): $p_{(1)}=0{,}95$. Caso (b): 200 tokens de probabilidad $0{,}005$.
\planteamiento \emph{top-k} fija el número de tokens; el núcleo fija la masa conservada (\cref{prop:nucleo}) y, por la \cref{prop:tv-truncada}, la masa eliminada es la distorsión.
\pasos
\begin{enumerate}[label=\arabic*.]
  \item (a) Núcleo: la acumulada del primer token es $0{,}95\ge 0{,}9$; se conserva \textbf{1} token y se elimina una masa $0{,}05$.
  \item (a) \emph{Top-k}: se conservan 50 tokens; los 49 últimos se reparten como mucho la masa $0{,}05$ y son cola de baja calidad.
  \item (b) Núcleo: hace falta $k^{\ast}$ con $0{,}005\,k^{\ast}\ge 0{,}9$, es decir, $k^{\ast}=0{,}9/0{,}005=180$ tokens; se elimina $1-0{,}9=0{,}1$.
  \item (b) \emph{Top-k}: los 50 tokens conservados suman $50\cdot 0{,}005=0{,}25$; se elimina una masa $0{,}75$ de opciones razonables.
\end{enumerate}
\resultado \textbf{(a) núcleo: 1 token; \emph{top-k}: 50. (b) núcleo: 180 tokens, elimina $0{,}1$; \emph{top-k}: 50 tokens, elimina $0{,}75$.} El núcleo adapta el tamaño del conjunto a la confianza del modelo.
\comprobacion en (b), $\TV=0{,}75$ para \emph{top-k} frente a $\TV=0{,}1\le 1-p$ para el núcleo, como garantiza la \cref{prop:tv-truncada}.
\end{solucion}

\begin{solucion}{eje:1-logits}
\datos $z=(1,0,-1)$ para $A,B,C$.
\planteamiento $p_\tau=\softmax(z/\tau)$ (\cref{def:temperatura}); \emph{top-k} y núcleo según la \cref{def:topk-topp}.
\pasos
\begin{enumerate}[label=\arabic*.]
  \item $\tau=1$: $e^{1}=2{,}718$, $e^{0}=1$, $e^{-1}=0{,}3679$; suma $4{,}086$. $p_1=(0{,}6652;\allowbreak\ 0{,}2447;\allowbreak\ 0{,}0900)$.
  \item $\tau=0{,}5$: $z/\tau=(2,0,-2)$; $7{,}389+1+0{,}1353=8{,}524$. $p_{0{,}5}=(0{,}8668;\allowbreak\ 0{,}1173;\allowbreak\ 0{,}0159)$.
  \item \emph{Top-2}: masa $0{,}6652+0{,}2447=0{,}9100$; distribución $(0{,}6652/0{,}9100;\allowbreak\ 0{,}2447/0{,}9100)=(0{,}7311;\allowbreak\ 0{,}2689)$.
  \item Núcleo $0{,}9$: acumuladas $0{,}6652<0{,}9$ y $0{,}9100\ge 0{,}9$; se conservan $A$ y $B$.
  \item La distribución \emph{top-2} es $e^{z_1}/(e^{z_1}+e^{z_2})=1/(1+e^{-(z_1-z_2)})$, que solo depende de $z_1-z_2=1$: $1/(1+0{,}3679)=0{,}7311$.
\end{enumerate}
\resultado \textbf{$p_1=(0{,}6652;\allowbreak\,0{,}2447;\allowbreak\,0{,}0900)$; $p_{0{,}5}=(0{,}8668;\allowbreak\,0{,}1173;\allowbreak\,0{,}0159)$; \emph{top-2}: $(0{,}7311;\allowbreak\,0{,}2689)$; el núcleo $0{,}9$ conserva $A$ y $B$.}
\comprobacion $0{,}6652+0{,}2447+0{,}0900=0{,}9999\approx 1$ y $0{,}8668+0{,}1173+0{,}0159=1$.
\end{solucion}

\begin{solucion}{eje:1-slp}
\datos $u=(1{,}2;\allowbreak\ 0{,}9;\allowbreak\ 0{,}1;\allowbreak\ -0{,}5)$; $\tau\in\{1;\allowbreak\ 0{,}2;\allowbreak\ 2\}$.
\planteamiento $p_\tau(i)=e^{u_i/\tau}/\sum_j e^{u_j/\tau}$ y la \cref{prop:limites-tau}.
\pasos
\begin{enumerate}[label=\arabic*.]
  \item $\tau=1$: $e^{u}=(3{,}320;\allowbreak\ 2{,}460;\allowbreak\ 1{,}105;\allowbreak\ 0{,}6065)$, suma $7{,}491$; $p_1=(0{,}4432;\allowbreak\ 0{,}3283;\allowbreak\ 0{,}1475;\allowbreak\ 0{,}0810)$.
  \item $\tau=0{,}2$: $u/\tau=(6;\allowbreak\ 4{,}5;\allowbreak\ 0{,}5;\allowbreak\ -2{,}5)$; $e^{u/\tau}=(403{,}4;\allowbreak\ 90{,}02;\allowbreak\ 1{,}649;\allowbreak\ 0{,}0821)$, suma $495{,}2$; $p_{0{,}2}=(0{,}8147;\allowbreak\ 0{,}1818;\allowbreak\ 0{,}0033;\allowbreak\ 0{,}0002)$.
  \item $\tau=2$: $u/\tau=(0{,}6;\allowbreak\ 0{,}45;\allowbreak\ 0{,}05;\allowbreak\ -0{,}25)$; $e^{u/\tau}=(1{,}822;\allowbreak\ 1{,}568;\allowbreak\ 1{,}051;\allowbreak\ 0{,}7788)$, suma $5{,}221$; $p_2=(0{,}3490;\allowbreak\ 0{,}3004;\allowbreak\ 0{,}2014;\allowbreak\ 0{,}1492)$.
  \item Sumas: $0{,}4432+0{,}3283+0{,}1475+0{,}0810=1{,}0000$; $0{,}8147+0{,}1818+0{,}0033+0{,}0002=1{,}0000$; $0{,}3490+0{,}3004+0{,}2014+0{,}1492=1{,}0000$.
  \item (b) El máximo $1{,}2$ es único: por la \cref{prop:limites-tau}, $p_\tau\to(1,0,0,0)$ si $\tau\to 0$ (búsqueda voraz) y $p_\tau\to(\frac14,\frac14,\frac14,\frac14)$ si $\tau\to\infty$.
\end{enumerate}
\resultado \textbf{$p_1=(0{,}4432;\allowbreak\,0{,}3283;\allowbreak\,0{,}1475;\allowbreak\,0{,}0810)$, $p_{0{,}2}=(0{,}8147;\allowbreak\,0{,}1818;\allowbreak\,0{,}0033;\allowbreak\,0{,}0002)$, $p_2=(0{,}3490;\allowbreak\,0{,}3004;\allowbreak\,0{,}2014;\allowbreak\,0{,}1492)$; voraz si $\tau\to 0$, uniforme si $\tau\to\infty$.}
\comprobacion la tendencia es la esperada: la probabilidad del primer token crece al bajar $\tau$ ($0{,}3490<0{,}4432<0{,}8147$), y $p_1$ coincide con los valores $0{,}44$, $0{,}33$, $0{,}15$ y $0{,}08$ de la figura~7.17 de~\cite{jm3}.
\end{solucion}

\begin{solucion}{eje:1-nucleo}
\datos $v=(1;\allowbreak\ 0{,}5;\allowbreak\ 0;\allowbreak\ -1)$; $\tau\in\{1;\allowbreak\ 0{,}5\}$; $p=0{,}8$.
\planteamiento softmax con temperatura y, después, núcleo: ordenar, acumular y cortar en cuanto la suma llega a $0{,}8$.
\pasos
\begin{enumerate}[label=\arabic*.]
  \item $\tau=1$: $e^{v}=(2{,}718;\allowbreak\ 1{,}649;\allowbreak\ 1;\allowbreak\ 0{,}3679)$, suma $5{,}735$; $p_1=(0{,}4740;\allowbreak\ 0{,}2875;\allowbreak\ 0{,}1744;\allowbreak\ 0{,}0641)$.
  \item Acumuladas: $0{,}4740$; $0{,}7615$; $0{,}9359$. Se conservan tres tokens; renormalizando por $0{,}9359$: $(0{,}5065;\allowbreak\ 0{,}3072;\allowbreak\ 0{,}1863;\allowbreak\ 0)$.
  \item $\tau=0{,}5$: $v/\tau=(2;\allowbreak\ 1;\allowbreak\ 0;\allowbreak\ -2)$, $e^{v/\tau}=(7{,}389;\allowbreak\ 2{,}718;\allowbreak\ 1;\allowbreak\ 0{,}1353)$, suma $11{,}24$; $p_{0{,}5}=(0{,}6572;\allowbreak\ 0{,}2418;\allowbreak\ 0{,}0889;\allowbreak\ 0{,}0120)$.
  \item Acumuladas: $0{,}6572$; $0{,}8990$. Bastan dos tokens; renormalizando por $0{,}8990$: $(0{,}7311;\allowbreak\ 0{,}2689;\allowbreak\ 0;\allowbreak\ 0)$.
  \item (c) Si $\tau\to 0$, la distribución tiende a $(1,0,0,0)$ (\cref{prop:limites-tau}) y el núcleo conserva un solo token: búsqueda voraz.
\end{enumerate}
\resultado \textbf{(a) $p_1=(0{,}4740;\allowbreak\,0{,}2875;\allowbreak\,0{,}1744;\allowbreak\,0{,}0641)$, $p_{0{,}5}=(0{,}6572;\allowbreak\,0{,}2418;\allowbreak\,0{,}0889;\allowbreak\,0{,}0120)$. (b) $\tau=1$: 3 tokens, $(0{,}5065;\allowbreak\,0{,}3072;\allowbreak\,0{,}1863;\allowbreak\,0)$; $\tau=0{,}5$: 2 tokens, $(0{,}7311;\allowbreak\,0{,}2689;\allowbreak\,0;\allowbreak\,0)$. (c) Voraz.}
\comprobacion con $\tau=0{,}5$ la diferencia entre los dos primeros logits escalados es $2-1=1$ y $1/(1+e^{-1})=0{,}7311$.
\end{solucion}

\begin{solucion}{eje:1-haz}
\datos $K=50\,000$, $T=20$, $m=4$; candidatas $(L=3,\ P=0{,}12)$ y $(L=5,\ P=0{,}08)$; $a=0{,}75$.
\planteamiento la voraz evalúa $K$ extensiones por paso y el haz $mK$ (\cref{prop:coste-haz}); la exhaustiva, $K^{T}$ (\cref{prop:exhaustiva}). Puntuación de la \cref{obs:longitud}.
\pasos
\begin{enumerate}[label=\arabic*.]
  \item Voraz: $KT=50\,000\cdot 20=10^{6}$.
  \item Haz: $mKT=4\cdot 10^{6}$.
  \item Exhaustiva: $\log_{10}K^{T}=20\cdot 4{,}699=93{,}98$, luego $K^{T}\approx 9{,}537\cdot 10^{93}$.
  \item Candidata 1: $\ln 0{,}12=-2{,}120$; $3^{0{,}75}=e^{0{,}75\cdot 1{,}099}=e^{0{,}8240}=2{,}280$; puntuación $-2{,}120/2{,}280=-0{,}9301$.
  \item Candidata 2: $\ln 0{,}08=-2{,}526$; $5^{0{,}75}=e^{0{,}75\cdot 1{,}609}=e^{1{,}207}=3{,}344$; puntuación $-2{,}526/3{,}344=-0{,}7554$.
  \item Con normalización gana la 2 ($-0{,}7554>-0{,}9301$); sin ella gana la 1 ($-2{,}120>-2{,}526$).
  \item (c) Sí: con $m\ge K^{T-1}$ el haz no poda y es la búsqueda exhaustiva; con $m=1$ es la voraz (\cref{prop:coste-haz}).
\end{enumerate}
\resultado \textbf{(a) $10^{6}$, $4\cdot 10^{6}$ y $\approx 9{,}537\cdot 10^{93}$ evaluaciones. (b) Con $a=0{,}75$ gana la candidata larga; sin normalizar, la corta. (c) Sí, con un haz que no poda.}
\comprobacion sin normalización cada token añade un logaritmo negativo y las secuencias largas siempre salen perjudicadas, como indica la \cref{obs:longitud}.
\end{solucion}

\begin{solucion}{eje:1-potencia}
\datos $P=(0{,}6;\allowbreak\ 0{,}3;\allowbreak\ 0{,}1)$; $Q\propto P^{a}$.
\planteamiento por la \cref{prop:potencia}, $Q$ es la distribución con temperatura $\tau=1/a$.
\pasos
\begin{enumerate}[label=\arabic*.]
  \item (a) Duplicar la temperatura: $\tau=2$, $a=\frac12$. Reducirla a la mitad: $\tau=\frac12$, $a=2$.
  \item $a=2$: $P^{2}=(0{,}36;\allowbreak\ 0{,}09;\allowbreak\ 0{,}01)$, suma $0{,}46$; $Q=(0{,}7826;\allowbreak\ 0{,}1957;\allowbreak\ 0{,}0217)$.
  \item $a=\frac12$: $\sqrt P=(0{,}7746;\allowbreak\ 0{,}5477;\allowbreak\ 0{,}3162)$, suma $1{,}639$; $Q=(0{,}4727;\allowbreak\ 0{,}3343;\allowbreak\ 0{,}1930)$.
  \item (c) $\tau\to 0$ equivale a $a\to\infty$: domina el mayor $P(i)$ y $Q\to(1,0,0)$. $\tau\to\infty$ equivale a $a\to 0$: $P(i)^{0}=1$ y $Q\to(\frac13,\frac13,\frac13)$.
\end{enumerate}
\resultado \textbf{(a) $a=\frac12$ duplica $\tau$ y $a=2$ la reduce a la mitad. (b) $a=2$: $(0{,}7826;\allowbreak\,0{,}1957;\allowbreak\,0{,}0217)$; $a=\frac12$: $(0{,}4727;\allowbreak\,0{,}3343;\allowbreak\,0{,}1930)$. (c) Voraz y uniforme.}
\comprobacion las dos distribuciones suman~1; con $a=2$ se afila ($0{,}6\to 0{,}7826$) y con $a=\frac12$ se aplana ($0{,}6\to 0{,}4727$).
\end{solucion}

\begin{solucion}{eje:1-entropia}
\datos $p_\tau(i)=e^{z_i/\tau}/Z(\tau)$ con $Z(\tau)=\sum_je^{z_j/\tau}$; en (c), $z=(1,0,-1)$ y $\tau=1$.
\planteamiento se trabaja con $\beta=1/\tau$ y se aplica la regla de la cadena, $\frac{\mathrm{d}H}{\mathrm{d}\tau}=\frac{\mathrm{d}H}{\mathrm{d}\beta}\,\frac{\mathrm{d}\beta}{\mathrm{d}\tau}$.
\pasos
\begin{enumerate}[label=\arabic*.]
  \item Con $\beta=1/\tau$, $\ln p_\tau(i)=\beta z_i-\ln Z$ y $H=-\sum_ip_\tau(i)(\beta z_i-\ln Z)=\ln Z-\beta\,\E[z]$, donde $\E[z]=\sum_i p_\tau(i)z_i$.
  \item $\dfrac{\mathrm{d}\ln Z}{\mathrm{d}\beta}=\dfrac{\sum_jz_je^{\beta z_j}}{Z}=\E[z]$.
  \item $\dfrac{\mathrm{d}\E[z]}{\mathrm{d}\beta}=\sum_iz_i\,p_\tau(i)\bigl(z_i-\E[z]\bigr)=\E[z^{2}]-\E[z]^{2}=\Var_{p_\tau}(z)$, usando $\frac{\mathrm{d}p_\tau(i)}{\mathrm{d}\beta}=p_\tau(i)(z_i-\E[z])$.
  \item $\dfrac{\mathrm{d}H}{\mathrm{d}\beta}=\E[z]-\E[z]-\beta\Var(z)=-\beta\Var(z)$.
  \item $\dfrac{\mathrm{d}\beta}{\mathrm{d}\tau}=-\dfrac{1}{\tau^{2}}$, luego $\dfrac{\mathrm{d}H}{\mathrm{d}\tau}=\beta\Var(z)\dfrac{1}{\tau^{2}}=\dfrac{\Var(z)}{\tau^{3}}\ge 0$.
  \item (b) Como $p_\tau(i)>0$ para todo $i$, $\Var_{p_\tau}(z)=0$ si y solo si todos los $z_i$ son iguales; entonces $p_\tau$ es uniforme para todo $\tau$ y $H=\ln K$ es constante.
  \item (c) $p_1=(0{,}6652;\allowbreak\ 0{,}2447;\allowbreak\ 0{,}0900)$; $\E[z]=0{,}6652-0{,}0900=0{,}5752$; $\E[z^{2}]=0{,}6652+0{,}0900=0{,}7553$; $\Var(z)=0{,}7553-0{,}5752^{2}=0{,}7553-0{,}3309=0{,}4244$.
\end{enumerate}
\resultado \textbf{$\mathrm{d}H/\mathrm{d}\tau=\Var_{p_\tau}(z)/\tau^{3}\ge 0$; es nula solo si todos los logits son iguales; en el caso (c) vale $0{,}4244$.} \hfill$\blacksquare$
\comprobacion una diferencia centrada, $(H(1{,}001)-H(0{,}999))/0{,}002$, da $0{,}4244$; además, los valores del \cref{ej:muestreo} cumplen $H(0{,}5)=0{,}455<H(1)=0{,}948<H(2)=1{,}245$ nats.
\end{solucion}

% ---------------------------------------------------------------------
\section{Tema 2}

\begin{solucion}{eje:2-fases}
\datos el mismo Transformer; \emph{prefill} con $n$ tokens a la vez y decodificación con uno por paso.
\planteamiento comparar la intensidad respecto a los pesos con la intensidad crítica $\Icrit$ (\cref{prop:intensidad-capas}).
\pasos
\begin{enumerate}[label=\arabic*.]
  \item \emph{Prefill}: cada peso de 2 bytes se lee una vez y se usa con los $n$ tokens, a 2 FLOP por token: $2n/2=n$ FLOP/byte.
  \item Decodificación: cada peso se usa con un solo token: $2/2=1$ FLOP/byte.
  \item Con $\Icrit=150$ (\cref{ej:cota}), un prompt de $n=1000$ tokens supera la intensidad crítica y queda limitado por el cómputo; la decodificación, con $1\ll 150$, por la memoria.
\end{enumerate}
\resultado \textbf{el \emph{prefill} reutiliza cada peso $n$ veces y queda limitado por el cómputo; la decodificación lo usa una vez y queda limitada por el ancho de banda.}
\comprobacion un prompt de un solo token tendría la misma intensidad que un paso de decodificación: las dos fases solo se distinguen por cuántos tokens comparten cada lectura de pesos.
\end{solucion}

\begin{solucion}{eje:2-vf}
\planteamiento contrastar con el \cref{teo:cache}, la \cref{prop:lote} y la \cref{prop:tam-cache}.
\pasos
\begin{enumerate}[label=(\alph*)]
  \item \textbf{Falsa.} Por el \cref{teo:cache}, la caché guarda exactamente lo que se recalcularía; solo cambia el coste (salvo diferencias de redondeo en coma flotante).
  \item \textbf{Falsa.} Por la \cref{prop:lote}, $R(B)=B/(t_0+Bt_1)$ crece menos que linealmente; en la \cref{tab:lote}, pasar de $B=32$ a $B=64$ lleva el rendimiento de $1346$ a $1578$ tokens/s, un factor $1{,}17$.
  \item \textbf{Verdadera.} La caché es proporcional a $\nkv$, que se divide por $32/8=4$; en la misma memoria caben cuatro veces más secuencias de la misma longitud.
\end{enumerate}
\resultado \textbf{(a) falsa; (b) falsa; (c) verdadera.}
\end{solucion}

\begin{solucion}{eje:2-kv}
\datos $L=24$, $\nkv=16$, $\dhd=64$, $s=2$; $S=2048$; $B=16$.
\planteamiento \cref{prop:tam-cache}: $c_{\mathrm{kv}}=2L\nkv\dhd s$; se multiplica por tokens y por secuencias.
\pasos
\begin{enumerate}[label=\arabic*.]
  \item $c_{\mathrm{kv}}=2\cdot 24\cdot 16\cdot 64\cdot 2=98\,304$~bytes; $98\,304/1024=96$~KiB.
  \item Por secuencia: $96\cdot 2048=196\,608$~KiB $=196\,608/1024=192$~MiB.
  \item Lote: $192\cdot 16=3072$~MiB $=3$~GiB.
\end{enumerate}
\resultado \textbf{96~KiB por token, 192~MiB por secuencia y 3~GiB para el lote.}
\comprobacion $98\,304\cdot 2048\cdot 16=3\,221\,225\,472$~bytes $=3\cdot 1024^{3}$.
\end{solucion}

\begin{solucion}{eje:2-cota}
\datos $N=1{,}3\cdot 10^{10}$; $\Bmem=10^{12}$~bytes/s; $s\in\{2;\allowbreak\,1;\allowbreak\,0{,}5\}$.
\planteamiento \cref{teo:cota-memoria} con $B=1$: $\text{tokens/s}\le\Bmem/(sN)$.
\pasos
\begin{enumerate}[label=\arabic*.]
  \item fp16: $sN=2{,}6\cdot 10^{10}$~bytes; cota $10^{12}/2{,}6\cdot 10^{10}=38{,}46$ tokens/s.
  \item int8: $sN=1{,}3\cdot 10^{10}$~bytes; cota $76{,}92$ tokens/s.
  \item int4: $sN=6{,}5\cdot 10^{9}$~bytes; cota $153{,}8$ tokens/s.
\end{enumerate}
\resultado \textbf{$38{,}46$, $76{,}92$ y $153{,}8$ tokens/s.}
\comprobacion $153{,}8/38{,}46=4{,}0=2/0{,}5$, el cociente de bytes por peso.
\end{solucion}

\begin{solucion}{eje:2-llama13}
\datos $N=1{,}3015\cdot 10^{10}$; bf16 ($s=2$); $L=40$, $\nkv=40$, $\dhd=128$; $S=1024$; $\Bmem=3{,}35\cdot 10^{12}$~bytes/s; memoria de 80~GB.
\planteamiento memoria $=sN+B\,S\,c_{\mathrm{kv}}$; latencia por paso $=$ memoria$/\Bmem$ (\cref{teo:cota-memoria}); rendimiento $=B/$latencia.
\pasos
\begin{enumerate}[label=\arabic*.]
  \item $c_{\mathrm{kv}}=2\cdot 40\cdot 40\cdot 128\cdot 2=819\,200$~bytes; por secuencia, $1024\cdot 819\,200=8{,}389\cdot 10^{8}$~bytes $=0{,}8389$~GB.
  \item Pesos: $2\cdot 1{,}3015\cdot 10^{10}=26{,}03$~GB.
  \item $B=1$: memoria $26{,}03+0{,}84=26{,}87$~GB; latencia $26{,}87\cdot 10^{9}/3{,}35\cdot 10^{12}=8{,}02$~ms; rendimiento $1/0{,}00802=124{,}7$ tokens/s.
  \item $B=64$: caché $64\cdot 0{,}8389=53{,}69$~GB; memoria $79{,}72$~GB; latencia $23{,}80$~ms; rendimiento $64/0{,}02380=2689$ tokens/s.
  \item $B=256$: caché $214{,}7$~GB; memoria $240{,}8$~GB; latencia $71{,}87$~ms; rendimiento $256/0{,}07187=3562$ tokens/s.
  \item (c) $B=1$ y $B=64$ caben en 80~GB (este último por muy poco); $B=256$ no.
\end{enumerate}
\resultado \textbf{caché de $0{,}839$~GB por secuencia; $(B,\text{memoria},\text{latencia},\text{rendimiento})=(1;\allowbreak\ 26{,}87~\text{GB};\allowbreak\ 8{,}02~\text{ms};\allowbreak\ 124{,}7)$, $(64;\allowbreak\ 79{,}72;\allowbreak\ 23{,}80;\allowbreak\ 2689)$, $(256;\allowbreak\ 240{,}8;\allowbreak\ 71{,}87;\allowbreak\ 3562)$.} Aumentar el lote empeora la latencia y mejora el rendimiento con rendimientos decrecientes; con $B=256$ la caché no cabe en la GPU.
\comprobacion el límite de la \cref{prop:lote} es $\Bmem/(S\,c_{\mathrm{kv}})=3{,}35\cdot 10^{12}/8{,}389\cdot 10^{8}=3993$ tokens/s, por encima de los tres valores y cada vez más próximo.
\end{solucion}

\begin{solucion}{eje:2-intensidad}
\datos lote $B$; pesos en fp16 (2 bytes); $\Icrit=150$~FLOP/byte.
\planteamiento contar, para un peso, las operaciones y los bytes en un paso.
\pasos
\begin{enumerate}[label=\arabic*.]
  \item En un paso, cada peso $w$ se lee una vez (2 bytes) y participa en una multiplicación y una suma por cada una de las $B$ secuencias: $2B$ FLOP.
  \item Intensidad respecto a los pesos: $2B/2=B$ FLOP/byte. \hfill$\blacksquare$
  \item Deja de estar limitado por los pesos si $B\ge\Icrit=150$.
  \item La caché KV es distinta para cada secuencia: cada número de la caché se lee en cada paso y se usa con una sola consulta, con intensidad menor que~1 (\cref{prop:intensidad-capas}(b)).
  \item Su tamaño crece con $B$: con Llama~2 7B y 2000 tokens, para $B=64$ ya son $67{,}1$~GB frente a $14$~GB de pesos (\cref{tab:lote}), y puede no caber en la memoria.
\end{enumerate}
\resultado \textbf{la intensidad es $B$ FLOP/byte; harían falta unas 150 secuencias por lote, pero la caché, con intensidad menor que~1 y tamaño proporcional a $B$, domina antes el tráfico y la memoria.}
\comprobacion con $B=1$ se recupera la intensidad de $1$~FLOP/byte del \cref{ej:cota}.
\end{solucion}

\begin{solucion}{eje:2-sincache}
\datos una capa de atención de dimensión $D$; prompt vacío; $T$ tokens generados.
\planteamiento el producto de $a$ consultas por $b$ claves cuesta $2abD$ FLOP, y el de los pesos de atención por los valores, otros $2abD$ (\cref{prop:intensidad-capas}).
\pasos
\begin{enumerate}[label=\arabic*.]
  \item Sin caché, el paso $t$ recalcula la atención completa sobre $t$ posiciones: $a=b=t$, coste $4Dt^{2}$.
  \item Total: $\sum_{t=1}^{T}4Dt^{2}=4D\frac{T(T+1)(2T+1)}{6}=\frac23DT(T+1)(2T+1)=\Theta(T^{3})$.
  \item Con caché, el paso $t$ calcula una consulta frente a $t$ claves: $a=1$, $b=t$, coste $4Dt$.
  \item Total: $\sum_{t=1}^{T}4Dt=4D\frac{T(T+1)}{2}=2DT(T+1)=\Theta(T^{2})$. \hfill$\blacksquare$
  \item Cociente: $\frac{\frac23DT(T+1)(2T+1)}{2DT(T+1)}=\frac{2T+1}{3}$; con $T=1000$, $2001/3=667$.
\end{enumerate}
\resultado \textbf{$\Theta(T^{3})$ sin caché frente a $\Theta(T^{2})$ con caché; para $T=1000$ la caché reduce el coste de la atención unas 667 veces.}
\comprobacion para $T=1$ ambos costes valen $4D$ y el cociente es $3/3=1$, como debe ser: con un solo token no hay nada que reutilizar.
\end{solucion}

% ---------------------------------------------------------------------
\section{Tema 3}

\begin{solucion}{eje:3-vf}
\planteamiento contrastar con los \cref{teo:especulativa,teo:especulativa-secuencias} y la \cref{def:ronda}.
\pasos
\begin{enumerate}[label=(\alph*)]
  \item \textbf{Falsa.} Por el \cref{teo:especulativa-secuencias}, la salida tiene exactamente la distribución del objetivo; solo cambia el tiempo.
  \item \textbf{Verdadera.} Si $q=p$, $\min(1,p/q)=1$ para todo token y $\beta=1$ (\cref{lem:residuo}).
  \item \textbf{Verdadera.} Si se rechaza la primera propuesta se emite un token del residuo; si no, al menos esa propuesta. Si se aceptan todas se emiten $\gamma+1$.
\end{enumerate}
\resultado \textbf{(a) falsa; (b) verdadera; (c) verdadera.}
\end{solucion}

\begin{solucion}{eje:3-dos}
\datos $p=(0{,}6;\allowbreak\ 0{,}4)$, $q=(0{,}9;\allowbreak\ 0{,}1)$.
\planteamiento aceptación $\min(1,p/q)$; $\beta=\sum\min(p,q)$; residuo $\max(0,p-q)$ normalizado; $P(\text{emitir }x)=\min(p,q)+(1-\beta)r(x)$.
\pasos
\begin{enumerate}[label=\arabic*.]
  \item Aceptar $A$: $\min(1;\allowbreak\ 0{,}6/0{,}9)=2/3\approx 0{,}6667$. Aceptar $B$: $\min(1;\allowbreak\ 0{,}4/0{,}1)=\min(1;4)=1$.
  \item $\beta=\min(0{,}6;0{,}9)+\min(0{,}4;0{,}1)=0{,}6+0{,}1=0{,}7$.
  \item Residuo: $\max(0;-0{,}3)=0$ y $\max(0;0{,}3)=0{,}3$; normalizado, $(0;\allowbreak\ 1)$.
  \item $P(A)=0{,}9\cdot\frac23=0{,}6$; $P(B)=0{,}1\cdot 1+0{,}3\cdot 1=0{,}4$.
\end{enumerate}
\resultado \textbf{aceptación $2/3$ para $A$ y $1$ para $B$; $\beta=0{,}7$; residuo $(0,1)$; salida $(0{,}6;\allowbreak\ 0{,}4)=p$.}
\comprobacion $1-\beta=0{,}3=\frac12(\abs{0{,}6-0{,}9}+\abs{0{,}4-0{,}1})$, como afirma la \cref{prop:beta-tv}.
\end{solucion}

\begin{solucion}{eje:3-cs336}
\datos $\V=\{A,B\}$; $q(A)>p(A)$; en (c), $p=(0{,}3;\allowbreak\ 0{,}7)$ y $q=(0{,}8;\allowbreak\ 0{,}2)$.
\planteamiento se aplica la regla de aceptación y rechazo token a token.
\pasos
\begin{enumerate}[label=\arabic*.]
  \item (a) $q(B)=1-q(A)<1-p(A)=p(B)$. Residuo: $\max(0,p(A)-q(A))=0$ y $\max(0,p(B)-q(B))=p(B)-q(B)>0$; normalizado, $(0,1)$.
  \item (b) Se propone $A$ con probabilidad $q(A)$ y se acepta con $p(A)/q(A)<1$: $P(\text{emitir }A)=q(A)\cdot\frac{p(A)}{q(A)}+0=p(A)$, porque el residuo nunca da $A$.
  \item Se emite $B$ si se propone y se acepta (probabilidad $q(B)\cdot 1$) o si se propone $A$, se rechaza y el residuo da $B$ (probabilidad $q(A)\bigl(1-\frac{p(A)}{q(A)}\bigr)\cdot 1=q(A)-p(A)$).
  \item $P(\text{emitir }B)=q(B)+q(A)-p(A)=1-p(A)=p(B)$.
  \item (c) Aceptar $A$: $0{,}3/0{,}8=0{,}375$; $P(A)=0{,}8\cdot 0{,}375=0{,}3$; $P(B)=0{,}2+0{,}8-0{,}3=0{,}7$.
\end{enumerate}
\resultado \textbf{$P(\text{emitir }A)=p(A)$ y $P(\text{emitir }B)=p(B)$; en (c), $(0{,}3;\allowbreak\ 0{,}7)$.}
\comprobacion $\beta=\min(0{,}3;0{,}8)+\min(0{,}7;0{,}2)=0{,}5$ y $\TV(p,q)=\frac12(0{,}5+0{,}5)=0{,}5=1-\beta$.
\end{solucion}

\begin{solucion}{eje:3-gamma}
\datos $\alpha=0{,}5$; $\gamma\in\{2,3,10\}$; $c=0{,}1$.
\planteamiento $\E[N_{\mathrm{r}}]=(1-\alpha^{\gamma+1})/(1-\alpha)$ (\cref{teo:tokens-ronda}) y $f(\gamma)=\E[N_{\mathrm{r}}]/(\gamma c+1)$ (\cref{teo:mejora}).
\pasos
\begin{enumerate}[label=\arabic*.]
  \item $\gamma=3$: $0{,}5^{4}=0{,}0625$; $\E[N_{\mathrm{r}}]=0{,}9375/0{,}5=1{,}875$.
  \item $\gamma=10$: $0{,}5^{11}=0{,}000488$; $\E[N_{\mathrm{r}}]=0{,}9995/0{,}5=1{,}999$.
  \item $\gamma=2$: $0{,}5^{3}=0{,}125$; $\E[N_{\mathrm{r}}]=0{,}875/0{,}5=1{,}75$.
  \item Factores: $f(2)=1{,}75/1{,}2=1{,}458$; $f(3)=1{,}875/1{,}3=1{,}442$; $f(10)=1{,}999/2=0{,}9995$.
\end{enumerate}
\resultado \textbf{$1{,}875$ y $1{,}999$ tokens por ronda; $f(2)=1{,}458$, $f(3)=1{,}442$ y $f(10)=0{,}9995$.} Con $\gamma=10$ se ganan $0{,}124$ tokens por ronda a cambio de 7 pasadas del borrador, y el resultado es más lento que no especular. Con $\alpha$ bajo conviene $\gamma$ pequeño; aquí el óptimo es $\gamma=2$.
\comprobacion el límite de $\E[N_{\mathrm{r}}]$ es $1/(1-\alpha)=2$, y con $\gamma=3$ ya se tiene el $93{,}75\,\%$ de ese máximo.
\end{solucion}

\begin{solucion}{eje:3-tv}
\datos $\tilde x\sim q$; $\beta=\sum_x\min(p(x),q(x))$; $p$ y $q$ son distribuciones.
\planteamiento (a) sin rechazo la salida es la propuesta; (b) escribir $1-\beta$ como suma de partes positivas.
\pasos
\begin{enumerate}[label=\arabic*.]
  \item (a) Si se acepta siempre, el token emitido es $\tilde x$, y $P(\text{emitir }x)=q(x)$, que solo coincide con $p$ si $p=q$.
  \item (b) Por el \cref{lem:residuo}, $1-\beta=\sum_x\max(0,p(x)-q(x))$.
  \item Como $\sum_x(p-q)=1-1=0$, las partes positiva y negativa suman lo mismo: $\sum_x\max(0,p-q)=\sum_x\max(0,q-p)$.
  \item Para cada $x$, $\abs{p-q}=\max(0,p-q)+\max(0,q-p)$; sumando, $\sum_x\abs{p-q}=2\sum_x\max(0,p-q)$.
  \item Luego $1-\beta=\frac12\sum_x\abs{p(x)-q(x)}=\TV(p,q)$. \hfill$\blacksquare$
\end{enumerate}
\resultado \textbf{sin rechazo la salida sigue $q$; la probabilidad de rechazo es $1-\beta=\TV(p,q)$.}
\comprobacion en el \cref{ej:especulativa}, $\frac12(0{,}3+0{,}3+0)=0{,}3=1-0{,}7$.
\end{solucion}

\begin{solucion}{eje:3-optimo}
\datos $0<\alpha<1$, $c>0$; $\E[N_{\mathrm{r}}](\gamma)=\sum_{j=0}^{\gamma}\alpha^{j}$; $f(\gamma)=\E[N_{\mathrm{r}}](\gamma)/(\gamma c+1)$.
\planteamiento acotar la suma geométrica y estudiar el cociente.
\pasos
\begin{enumerate}[label=\arabic*.]
  \item (a) $\E[N_{\mathrm{r}}](\gamma+1)-\E[N_{\mathrm{r}}](\gamma)=\alpha^{\gamma+1}>0$, luego es estrictamente creciente; y $\sum_{j=0}^{\gamma}\alpha^{j}<\sum_{j=0}^{\infty}\alpha^{j}=1/(1-\alpha)$.
  \item (b) $0<f(\gamma)<\dfrac{1}{(1-\alpha)(\gamma c+1)}\to 0$ cuando $\gamma\to\infty$.
  \item Sea $f_1=f(1)>0$. Existe $\gamma_0$ tal que $f(\gamma)<f_1$ para $\gamma>\gamma_0$; el máximo de $f$ sobre $\{1,\dots,\gamma_0\}$, que existe por ser un conjunto finito, es el máximo global.
  \item (c) Como $\alpha^{j}\le\alpha$ para $j\ge 1$, $\E[N_{\mathrm{r}}](\gamma)\le 1+\gamma\alpha\le 1+\gamma c$ si $\alpha\le c$, y entonces $f(\gamma)\le 1$. \hfill$\blacksquare$
\end{enumerate}
\resultado \textbf{$\E[N_{\mathrm{r}}]$ crece estrictamente y está acotado por $1/(1-\alpha)$; $f(\gamma)\to 0$ y hay un $\gamma$ óptimo finito; si $\alpha\le c$, especular nunca compensa.} Junto con el \cref{cor:alpha-c}, la condición $\alpha>c$ es necesaria y suficiente para que exista algún $\gamma$ con $f(\gamma)>1$.
\comprobacion con $\alpha=0{,}5$ y $c=0{,}1$ (\cref{eje:3-gamma}), $f$ vale $1{,}364$; $1{,}458$; $1{,}442$ para $\gamma=1,2,3$ y $0{,}9995$ para $\gamma=10$: sube, alcanza el máximo en $\gamma=2$ y decrece.
\end{solucion}

% ---------------------------------------------------------------------
\section{Tema 4}

\begin{solucion}{eje:4-conceptos}
\planteamiento relacionar cada pregunta con las \cref{def:afin,def:simetrico,def:granularidad} y el \cref{teo:error-cuant}.
\pasos
\begin{enumerate}[label=(\alph*)]
  \item Los pesos se distribuyen de forma aproximadamente simétrica alrededor de~0; el cuantizador simétrico aprovecha bien sus niveles y evita el término de punto cero en el producto (\cref{prop:producto-entero}). Una activación no negativa, como la salida de una ReLU, usaría solo la mitad de los niveles simétricos; el afín coloca los $2^{b}$ niveles en $[0,x_{\max}]$ y reduce el paso a la mitad.
  \item Los ceros son muy frecuentes (relleno, activaciones nulas, pesos podados); si el~0 se representa exactamente, no añaden error ni sesgo.
  \item La escala por canal se adapta al rango de cada canal: un canal de valores pequeños no hereda el paso grande impuesto por otro con valores grandes (\cref{prop:atipico}), y la descuantización sigue siendo un reescalado del producto entero (\cref{prop:factorizar}).
\end{enumerate}
\resultado \textbf{simétrica para pesos centrados y afín para rangos asimétricos; el~0 exacto evita errores en valores muy frecuentes; la escala por canal reduce el error sin romper la aritmética entera.}
\end{solucion}

\begin{solucion}{eje:4-int4}
\datos $x=(-1;\allowbreak\ 0;\allowbreak\ 0{,}62;\allowbreak\ 2)$; $b=4$, enteros $0,\dots,15$.
\planteamiento \cref{def:afin} con calibración por mínimo y máximo.
\pasos
\begin{enumerate}[label=\arabic*.]
  \item $\Delta=(2-(-1))/15=3/15=0{,}2$.
  \item $Z=\rnd(1/0{,}2)=\rnd(5)=5$.
  \item $x/\Delta=(-5;\allowbreak\ 0;\allowbreak\ 3{,}1;\allowbreak\ 10)$; redondeo: $(-5;\allowbreak\ 0;\allowbreak\ 3;\allowbreak\ 10)$.
  \item $Q=(0;\allowbreak\ 5;\allowbreak\ 8;\allowbreak\ 15)$; todos en $[0,15]$, sin recorte.
  \item $\hat x=0{,}2\cdot(-5;\allowbreak\ 0;\allowbreak\ 3;\allowbreak\ 10)=(-1;\allowbreak\ 0;\allowbreak\ 0{,}6;\allowbreak\ 2)$; errores $(0;\allowbreak\ 0;\allowbreak\ 0{,}02;\allowbreak\ 0)$.
  \item $x=3$: $\rnd(3/0{,}2)+5=15+5=20$, recortado a $15$; $\hat x=0{,}2\cdot 10=2$ y el error es $1$.
\end{enumerate}
\resultado \textbf{$\Delta=0{,}2$; $Z=5$; $Q=(0;\allowbreak\,5;\allowbreak\,8;\allowbreak\,15)$; $\hat x=(-1;\allowbreak\,0;\allowbreak\,0{,}6;\allowbreak\,2)$; error máximo $0{,}02$. Con $x=3$ hay saturación y el error es $1$.}
\comprobacion $0{,}02\le\Delta/2=0{,}1$; el~0 se reconstruye exactamente y el error de saturación, $1>0{,}1$, no respeta la cota, como prevé el \cref{teo:error-cuant}.
\end{solucion}

\begin{solucion}{eje:4-3bits}
\datos la matriz $W$; $b=3$, enteros $0,\dots,7$; $x_{\min}=-1{,}6$, $x_{\max}=2{,}6$.
\planteamiento \cref{def:afin} por tensor.
\pasos
\begin{enumerate}[label=\arabic*.]
  \item $\Delta=(2{,}6+1{,}6)/7=4{,}2/7=0{,}6$; $Z=\rnd(1{,}6/0{,}6)=\rnd(2{,}667)=3$.
  \item $W/\Delta$ y su redondeo:
  \[
  \frac{W}{\Delta}=\begin{pmatrix}2 & -0{,}6667 & 0{,}4167\\ -1{,}667 & 4{,}333 & 0\\ 0{,}8333 & -2{,}667 & 3{,}167\end{pmatrix},\qquad
  \rnd\Bigl(\frac{W}{\Delta}\Bigr)=\begin{pmatrix}2 & -1 & 0\\ -2 & 4 & 0\\ 1 & -3 & 3\end{pmatrix}.
  \]
  \item Sumando $Z=3$ (sin recorte) y reconstruyendo con $\hat W=0{,}6\,\rnd(W/\Delta)$:
  \[
  Q=\begin{pmatrix}5 & 2 & 3\\ 1 & 7 & 3\\ 4 & 0 & 6\end{pmatrix},\qquad
  \hat W=\begin{pmatrix}1{,}2 & -0{,}6 & 0\\ -1{,}2 & 2{,}4 & 0\\ 0{,}6 & -1{,}8 & 1{,}8\end{pmatrix}.
  \]
  \item Errores:
  \[
  W-\hat W=\begin{pmatrix}0 & 0{,}2 & 0{,}25\\ 0{,}2 & 0{,}2 & 0\\ -0{,}1 & 0{,}2 & 0{,}1\end{pmatrix}.
  \]
\end{enumerate}
\resultado \textbf{$\Delta=0{,}6$; $Z=3$; error máximo $0{,}25$, en el elemento $0{,}25$, que se reconstruye como~0.}
\comprobacion $0{,}25\le\Delta/2=0{,}3$. El extremo $-1{,}6$ no se reconstruye exactamente ($-1{,}8$) porque $Z$ se redondeó de $2{,}667$ a $3$, pero su error, $0{,}2$, respeta la cota, como garantiza el \cref{teo:error-cuant} también con $Z$ redondeado.
\end{solucion}

\begin{solucion}{eje:4-atipico}
\datos $v=(0{,}12;\allowbreak\,-0{,}31;\allowbreak\,0{,}05;\allowbreak\,0{,}27;\allowbreak\,-0{,}08;\allowbreak\,0{,}19;\allowbreak\,-0{,}22;\allowbreak\,6{,}4)$; int4 simétrico, $\Delta=\max\abs{v}/7$.
\planteamiento \cref{def:simetrico} con distintos grupos.
\pasos
\begin{enumerate}[label=\arabic*.]
  \item (a) Por tensor: $\Delta=6{,}4/7=0{,}9143$. Todos los valores salvo el atípico cumplen $\abs{v_i}<\Delta/2=0{,}457$ y se cuantizan a~0 (\cref{prop:atipico}); el atípico da $Q=7$ y $\hat v=6{,}4$. Error máximo: $0{,}31$.
  \item (b) Grupo 1, $(0{,}12;\allowbreak\,-0{,}31;\allowbreak\,0{,}05;\allowbreak\,0{,}27)$: $\Delta=0{,}31/7=0{,}04429$; $v/\Delta=(2{,}71;\allowbreak\,-7;\allowbreak\,1{,}13;\allowbreak\,6{,}10)$, $Q=(3;\allowbreak\,-7;\allowbreak\,1;\allowbreak\,6)$, $\hat v=(0{,}1329;\allowbreak\,-0{,}31;\allowbreak\,0{,}0443;\allowbreak\,0{,}2657)$; errores $(0{,}0129;\allowbreak\,0;\allowbreak\,0{,}0057;\allowbreak\,0{,}0043)$.
  \item Grupo 2, $(-0{,}08;\allowbreak\,0{,}19;\allowbreak\,-0{,}22;\allowbreak\,6{,}4)$: $\Delta=0{,}9143$; los tres primeros se anulan; errores $(0{,}08;\allowbreak\,0{,}19;\allowbreak\,0{,}22;\allowbreak\,0)$. Error máximo global: $0{,}22$.
  \item (c) Sin el atípico, $\Delta=0{,}31/7=0{,}04429$ para los 7 valores; $Q=(3;\allowbreak\,-7;\allowbreak\,1;\allowbreak\,6;\allowbreak\,-2;\allowbreak\,4;\allowbreak\,-5)$; $\hat v=(0{,}1329;\allowbreak\,-0{,}31;\allowbreak\,0{,}0443;\allowbreak\,0{,}2657;\allowbreak\,-0{,}0886;\allowbreak\,0{,}1771;\allowbreak\,-0{,}2214)$; error máximo $0{,}0129$, y el atípico se guarda exacto en fp16.
\end{enumerate}
\resultado \textbf{errores máximos: $0{,}31$ por tensor, $0{,}22$ por grupos de~4 y $0{,}0129$ separando el atípico.} Los grupos solo ayudan a los que no comparten grupo con el atípico; separarlo, como hace LLM.int8(), resuelve el problema.
\comprobacion en (c), $0{,}0129\le\Delta/2=0{,}0221$; en (a), la cota $\Delta/2=0{,}457$ es $20{,}6$ veces mayor que la de (c), igual que el cociente de escalas $6{,}4/0{,}31=20{,}6$.
\end{solucion}

\begin{solucion}{eje:4-entero}
\datos enteros en $[-127,127]$; acumulador hasta $2^{31}-1=2\,147\,483\,647$; escalas fp16 (16 bits); grupos de $64$ y $128$; $N=7\cdot 10^{9}$.
\planteamiento cada producto vale como mucho $127^{2}$ en módulo; \cref{prop:bits-grupo} para los bits.
\pasos
\begin{enumerate}[label=\arabic*.]
  \item $127^{2}=16\,129$. La suma de $n$ productos cumple $\abs{\sum}\le 16\,129\,n$, que no desborda si $16\,129\,n\le 2\,147\,483\,647$, es decir, $n\le 133\,144{,}25$.
  \item Mayor $n$: $133\,144$.
  \item Grupos de 64: $4+16/64=4{,}25$ bits por peso; memoria $7\cdot 10^{9}\cdot 4{,}25/8=3{,}719\cdot 10^{9}$~bytes.
  \item Grupos de 128: $4+16/128=4{,}125$ bits; memoria $7\cdot 10^{9}\cdot 4{,}125/8=3{,}609\cdot 10^{9}$~bytes.
\end{enumerate}
\resultado \textbf{(a) $n\le 133\,144$. (b) $4{,}25$ bits y $3{,}72$~GB con grupos de 64; $4{,}125$ bits y $3{,}61$~GB con grupos de 128.}
\comprobacion las dimensiones internas de los LLM (unos miles) quedan muy por debajo de $133\,144$, y $3{,}61$~GB coincide con la \cref{tab:formatos}.
\end{solucion}

\begin{solucion}{eje:4-ruido}
\datos $e\sim\mathcal{U}[-\Delta/2,\Delta/2]$; señal $x\sim\mathcal{U}[-A,A]$; $\Delta=2A/2^{b}$.
\planteamiento calcular los momentos con la densidad uniforme $1/\Delta$ y dividir varianzas.
\pasos
\begin{enumerate}[label=\arabic*.]
  \item $\E[e]=\frac{1}{\Delta}\int_{-\Delta/2}^{\Delta/2}e\,\mathrm{d}e=0$, por simetría del integrando impar.
  \item $\E[e^{2}]=\frac{1}{\Delta}\Bigl[\frac{e^{3}}{3}\Bigr]_{-\Delta/2}^{\Delta/2}=\frac{1}{\Delta}\cdot\frac{2}{3}\cdot\frac{\Delta^{3}}{8}=\frac{\Delta^{2}}{12}$. \hfill$\blacksquare$
  \item Varianza de la señal: con el mismo cálculo en $[-A,A]$, $\E[x^{2}]=\frac{(2A)^{2}}{12}=\frac{A^{2}}{3}$.
  \item Relación señal-ruido: $\dfrac{A^{2}/3}{\Delta^{2}/12}=\dfrac{4A^{2}}{\Delta^{2}}=\dfrac{4A^{2}}{4A^{2}/2^{2b}}=2^{2b}$.
  \item En decibelios: $10\log_{10}2^{2b}=20\,b\log_{10}2=6{,}0206\,b$. \hfill$\blacksquare$
  \item $b=4$: $24{,}08$~dB; $b=8$: $48{,}16$~dB.
\end{enumerate}
\resultado \textbf{$\E[e]=0$, $\E[e^{2}]=\Delta^{2}/12$; relación señal-ruido $2^{2b}$, es decir, $6{,}02\,b$~dB: $24{,}08$~dB con 4 bits y $48{,}16$~dB con 8.}
\comprobacion la simulación del texto ($\Delta=0{,}01$) da $8{,}345\cdot 10^{-6}$ frente a $8{,}333\cdot 10^{-6}$, y la diferencia entre 8 y 4 bits es $4\cdot 6{,}02=24{,}08$~dB, conforme al \cref{cor:6db}.
\end{solucion}

\begin{solucion}{eje:4-granularidad}
\datos cuantización simétrica; en (b), $X=(1\ 4)$, $W=(1\ 1)\T$, $\Delta^{X}=(1,4)$, $\Delta_W=1$.
\planteamiento comparar los máximos que definen las escalas; en (b), evaluar el sumatorio con y sin las escalas dentro.
\pasos
\begin{enumerate}[label=\arabic*.]
  \item (a) Por tensor, $\Delta=\max_{k,j}\abs{W_{kj}}/(2^{b-1}-1)$; por canal, $\Delta_j=\max_k\abs{W_{kj}}/(2^{b-1}-1)$.
  \item Para todo $j$, $\max_k\abs{W_{kj}}\le\max_{k,j'}\abs{W_{kj'}}$, luego $\Delta_j\le\Delta$ y $\Delta_j/2\le\Delta/2$. \hfill$\blacksquare$
  \item (b) $Q_X=(1/1,\ 4/4)=(1,1)$ y $Q_W=(1,1)\T$. Con las escalas dentro del sumatorio: $1\cdot 1\cdot 1\cdot 1+4\cdot 1\cdot 1\cdot 1=5=XW$.
  \item Producto entero: $Q_XQ_W=1+1=2$. Multiplicado por $\Delta_W=1$ y por una escala de $X$: $2\cdot 1=2$ o $2\cdot 4=8$; ninguno es $5$.
\end{enumerate}
\resultado \textbf{(a) la cota por canal nunca supera la cota por tensor. (b) El resultado correcto, 5, exige ponderar cada término por su escala $\Delta^{X}_{k}$; un reescalado único del producto entero da 2 u 8.}
\comprobacion la igualdad en (a) se da en el canal que contiene el máximo global, que es el único que no mejora.
\end{solucion}

% ---------------------------------------------------------------------
\section{Tema 5}

\begin{solucion}{eje:5-conceptos}
\planteamiento relacionar con la \cref{prop:atipico}, la \cref{prop:smooth} y la \cref{def:capas}.
\pasos
\begin{enumerate}[label=(\alph*)]
  \item Con absmax por token, $\Delta=60/127=0{,}472$; todos los valores de la fila con $\abs{x}<0{,}236$ se cuantizan a~0 y los demás usan muy pocos niveles (\cref{prop:atipico}). Como los atípicos aparecen en casi todas las filas y en las mismas dimensiones, el problema afecta a toda la matriz.
  \item Por la \cref{prop:smooth}, $XW=(X\diag(s)^{-1})(\diag(s)W)$ exactamente. $\diag(s)W$ se calcula una vez antes de servir el modelo, y $\diag(s)^{-1}$ se incorpora a los parámetros de la capa anterior, de modo que no hay operaciones nuevas en inferencia.
  \item GPTQ usa la matriz $H=2XX\T$, construida con datos de calibración, que describe cómo se combinan los errores de los pesos en la salida; el redondeo al más próximo trata cada peso por separado y no corrige los demás.
\end{enumerate}
\resultado \textbf{(a) el atípico infla la escala de toda la fila; (b) es una identidad que se precalcula; (c) la información de segundo orden de las entradas.}
\end{solucion}

\begin{solucion}{eje:5-absmax}
\datos $x=(3;\allowbreak\ -8;\allowbreak\ 5)$, $w=(0{,}6;\allowbreak\ 2{,}2;\allowbreak\ -4)$; enteros en $[-127,127]$.
\planteamiento $s_x=127/\norm{x}_\infty$, $s_w=127/\norm{w}_\infty$; $Q=\rnd(s\cdot\text{valor})$; $x\cdot w\approx(Q_x\cdot Q_w)/(s_xs_w)$ (\cref{def:simetrico}, \cref{prop:factorizar}).
\pasos
\begin{enumerate}[label=\arabic*.]
  \item $s_x=127/8=15{,}875$; $s_w=127/4=31{,}75$.
  \item $s_xx=(47{,}63;\allowbreak\ -127;\allowbreak\ 79{,}38)$, $Q_x=(48;\allowbreak\ -127;\allowbreak\ 79)$.
  \item $s_ww=(19{,}05;\allowbreak\ 69{,}85;\allowbreak\ -127)$, $Q_w=(19;\allowbreak\ 70;\allowbreak\ -127)$.
  \item Productos: $48\cdot 19=912$; $-127\cdot 70=-8890$; $79\cdot(-127)=-10\,033$. Suma en int32: $-18\,011$.
  \item Escala conjunta: $15{,}875\cdot 31{,}75=504{,}0$. Descuantizado: $-18\,011/504{,}0=-35{,}73$.
  \item Exacto: $1{,}8-17{,}6-20=-35{,}8$.
\end{enumerate}
\resultado \textbf{$x\cdot w\approx-35{,}73$ en int8 frente a $-35{,}8$ exacto: error relativo $0{,}07/35{,}8\approx 0{,}19\,\%$.}
\comprobacion los elementos de mayor módulo se cuantizan a $\pm 127$ sin error; los demás tienen error menor que medio paso, $0{,}5/15{,}875=0{,}0315$ en $x$ y $0{,}5/31{,}75=0{,}0157$ en $w$.
\end{solucion}

\begin{solucion}{eje:5-smooth}
\datos las matrices $X$ ($2\times 3$) y $W$ ($3\times 2$); $\alpha=\frac12$.
\planteamiento \cref{def:smooth}: $s_j=\sqrt{m^{X}_{j}/m^{W}_{j}}$; después $\hat X=X\diag(s)^{-1}$ y $\hat W=\diag(s)W$.
\pasos
\begin{enumerate}[label=\arabic*.]
  \item Máximos por canal: $m^{X}=(1;\allowbreak\ 16;\allowbreak\ 1)$ (columnas de $X$) y $m^{W}=(0{,}25;\allowbreak\ 0{,}25;\allowbreak\ 1)$ (filas de $W$).
  \item $s=(\sqrt{1/0{,}25};\allowbreak\ \sqrt{16/0{,}25};\allowbreak\ \sqrt{1/1})=(2;\allowbreak\ 8;\allowbreak\ 1)$.
  \item Dividiendo las columnas de $X$ por $s$ y multiplicando las filas de $W$ por $s$:
  \[
  \hat X=\begin{pmatrix}0{,}5 & -2 & 0{,}5\\ -0{,}25 & 1 & 1\end{pmatrix},\qquad
  \hat W=\begin{pmatrix}0{,}5 & -0{,}2\\ 0{,}4 & 2\\ -0{,}5 & 1\end{pmatrix}.
  \]
  \item Máximos tras el suavizado: $(0{,}5;\allowbreak\ 2;\allowbreak\ 1)$ en ambos, igual a $\sqrt{m^{X}_{j}m^{W}_{j}}=(\sqrt{0{,}25};\allowbreak\ \sqrt{4};\allowbreak\ \sqrt{1})$, como afirma la \cref{prop:alfa-medio}.
  \item $XW$: primera fila $(0{,}25-0{,}8-0{,}25;\allowbreak\ -0{,}1-4+0{,}5)=(-0{,}8;\allowbreak\ -3{,}6)$; segunda $(-0{,}125+0{,}4-0{,}5;\allowbreak\ 0{,}05+2+1)=(-0{,}225;\allowbreak\ 3{,}05)$.
  \item $\hat X\hat W$: primera fila $(0{,}25-0{,}8-0{,}25;\allowbreak\ -0{,}1-4+0{,}5)=(-0{,}8;\allowbreak\ -3{,}6)$; segunda $(-0{,}125+0{,}4-0{,}5;\allowbreak\ 0{,}05+2+1)=(-0{,}225;\allowbreak\ 3{,}05)$.
  \item Cociente entre máximos por canal de las activaciones: antes $16/1=16$; después $2/0{,}5=4$.
\end{enumerate}
\resultado \textbf{$s=(2;\allowbreak\,8;\allowbreak\,1)$; máximos por canal $(0{,}5;\allowbreak\,2;\allowbreak\,1)$ en $\hat X$ y en $\hat W$; $\hat X\hat W=XW=\begin{pmatrix}-0{,}8 & -3{,}6\\ -0{,}225 & 3{,}05\end{pmatrix}$; el desequilibrio entre canales de las activaciones baja de 16 a 4.}
\comprobacion el canal atípico (el segundo) se divide por $s_2=8$ en $X$ y su fila de $W$ se multiplica por~8: el producto de ambos máximos, $16\cdot 0{,}25=4=2\cdot 2$, se conserva.
\end{solucion}

\begin{solucion}{eje:5-gptq}
\datos $w=(0{,}7;\allowbreak\ 0{,}3)$; rejilla de paso $0{,}5$; $X=\begin{pmatrix}1 & 1\\ -1 & 0\end{pmatrix}$.
\planteamiento error $E(\delta)=\norm{\delta X}^{2}$; actualización del \cref{teo:obs} tras cuantizar $w_1$.
\pasos
\begin{enumerate}[label=\arabic*.]
  \item $XX\T=\begin{pmatrix}2 & -1\\ -1 & 1\end{pmatrix}$, $H=\begin{pmatrix}4 & -2\\ -2 & 2\end{pmatrix}$, $\det H=8-4=4$, $H^{-1}=\frac14\begin{pmatrix}2 & 2\\ 2 & 4\end{pmatrix}=\begin{pmatrix}0{,}5 & 0{,}5\\ 0{,}5 & 1\end{pmatrix}$.
  \item (b) Redondeo: $\quant(0{,}7)=0{,}5$, $\quant(0{,}3)=0{,}5$; $\delta=(-0{,}2;\allowbreak\ 0{,}2)$; $\delta X=(-0{,}2-0{,}2;\allowbreak\ -0{,}2+0)=(-0{,}4;\allowbreak\ -0{,}2)$; $E=0{,}16+0{,}04=0{,}20$.
  \item (c) Se cuantiza $w_1$: $\varepsilon=0{,}5-0{,}7=-0{,}2$; $\delta=\varepsilon\,(H^{-1})_{1,:}/[H^{-1}]_{11}=-0{,}2\cdot(0{,}5;\allowbreak\ 0{,}5)/0{,}5=(-0{,}2;\allowbreak\ -0{,}2)$.
  \item Pesos actualizados: $(0{,}5;\allowbreak\ 0{,}1)$, con error intermedio $\varepsilon^{2}/(2[H^{-1}]_{11})=0{,}04/1=0{,}04$.
  \item Se cuantiza $w_2=0{,}1$: $\quant(0{,}1)=0$. Resultado $\hat w=(0{,}5;\allowbreak\ 0)$, $\delta=(-0{,}2;\allowbreak\ -0{,}3)$.
  \item $\delta X=(-0{,}2+0{,}3;\allowbreak\ -0{,}2)=(0{,}1;\allowbreak\ -0{,}2)$; $E=0{,}01+0{,}04=0{,}05$.
\end{enumerate}
\resultado \textbf{$H^{-1}=\begin{pmatrix}0{,}5 & 0{,}5\\ 0{,}5 & 1\end{pmatrix}$; redondeo: $\hat w=(0{,}5;\allowbreak\,0{,}5)$, error $0{,}20$; GPTQ: $\hat w=(0{,}5;\allowbreak\,0)$, error $0{,}05$, cuatro veces menor.}
\comprobacion tras el primer paso, $\delta=(-0{,}2;-0{,}2)$ da $\delta X=(0;-0{,}2)$ y $E=0{,}04$, que coincide con la fórmula del \cref{teo:obs}. El peso $w_2$ se mueve en sentido contrario al redondeo ingenuo porque compensa el error de $w_1$ en la primera entrada.
\end{solucion}

\begin{solucion}{eje:5-qlora}
\datos $N=6{,}5\cdot 10^{10}$; fp16 (16 bits); NF4 (4 bits) con constantes de $0{,}5$ o $0{,}127$ bits por peso (\cref{prop:doble}).
\planteamiento memoria $=N\cdot(\text{bits por peso})/8$ bytes.
\pasos
\begin{enumerate}[label=\arabic*.]
  \item fp16: $6{,}5\cdot 10^{10}\cdot 2=1{,}3\cdot 10^{11}$~bytes $=130$~GB.
  \item NF4 con constantes fp32: $4{,}5$ bits; $6{,}5\cdot 10^{10}\cdot 4{,}5/8=3{,}656\cdot 10^{10}$~bytes $=36{,}56$~GB.
  \item Doble cuantización: $4+0{,}126953=4{,}127$ bits; $6{,}5\cdot 10^{10}\cdot 4{,}127/8=3{,}353\cdot 10^{10}$~bytes $=33{,}53$~GB.
  \item Ahorro: $36{,}56-33{,}53=3{,}03$~GB, que es $6{,}5\cdot 10^{10}\cdot 0{,}373/8$~bytes.
\end{enumerate}
\resultado \textbf{$130$~GB en fp16, $36{,}56$~GB en NF4 y $33{,}53$~GB con doble cuantización; la doble cuantización ahorra unos $3{,}03$~GB.}
\comprobacion $130/33{,}53=3{,}88$, algo menos que el factor $16/4=4$ por el coste de las constantes. Los 33~GB de pesos dejan margen en una GPU de 48~GB para los adaptadores y las activaciones.
\end{solucion}

\begin{solucion}{eje:5-destila}
\datos $v=(3,1,0)$; $z=(2,1,1)$; $\tau\in\{1,4\}$.
\planteamiento $p=\softmax(v/\tau)$, $q=\softmax(z/\tau)$ y $\partial C/\partial z=(q-p)/\tau$ (\cref{prop:grad-destilacion}).
\pasos
\begin{enumerate}[label=\arabic*.]
  \item $\tau=1$: $e^{v}=(20{,}09;\allowbreak\ 2{,}718;\allowbreak\ 1)$, suma $23{,}80$, $p=(0{,}8438;\allowbreak\ 0{,}1142;\allowbreak\ 0{,}0420)$; $e^{z}=(7{,}389;\allowbreak\ 2{,}718;\allowbreak\ 2{,}718)$, suma $12{,}83$, $q=(0{,}5761;\allowbreak\ 0{,}2119;\allowbreak\ 0{,}2119)$.
  \item Gradiente con $\tau=1$: $q-p=(-0{,}2677;\allowbreak\ 0{,}0977;\allowbreak\ 0{,}1699)$.
  \item $\tau=4$: $v/4=(0{,}75;\allowbreak\ 0{,}25;\allowbreak\ 0)$, $e^{v/4}=(2{,}117;\allowbreak\ 1{,}284;\allowbreak\ 1)$, suma $4{,}401$, $p=(0{,}4810;\allowbreak\ 0{,}2918;\allowbreak\ 0{,}2272)$; $z/4=(0{,}5;\allowbreak\ 0{,}25;\allowbreak\ 0{,}25)$, $e^{z/4}=(1{,}649;\allowbreak\ 1{,}284;\allowbreak\ 1{,}284)$, suma $4{,}217$, $q=(0{,}3910;\allowbreak\ 0{,}3045;\allowbreak\ 0{,}3045)$.
  \item Gradiente con $\tau=4$: $(q-p)/4=(-0{,}0900;\allowbreak\ 0{,}0127;\allowbreak\ 0{,}0773)/4=(-0{,}0225;\allowbreak\ 0{,}0032;\allowbreak\ 0{,}0193)$.
  \item (c) $\tau^{2}\cdot$ gradiente $=16\cdot(-0{,}0225;\allowbreak\ 0{,}0032;\allowbreak\ 0{,}0193)=(-0{,}360;\allowbreak\ 0{,}051;\allowbreak\ 0{,}309)$.
\end{enumerate}
\resultado \textbf{$\tau=1$: $p=(0{,}8438;\allowbreak\,0{,}1142;\allowbreak\,0{,}0420)$, $q=(0{,}5761;\allowbreak\,0{,}2119;\allowbreak\,0{,}2119)$, gradiente $(-0{,}2677;\allowbreak\,0{,}0977;\allowbreak\,0{,}1699)$. $\tau=4$: $p=(0{,}4810;\allowbreak\,0{,}2918;\allowbreak\,0{,}2272)$, $q=(0{,}3910;\allowbreak\,0{,}3045;\allowbreak\,0{,}3045)$, gradiente $(-0{,}0225;\allowbreak\,0{,}0032;\allowbreak\,0{,}0193)$.} Multiplicado por $\tau^{2}$ recupera un orden de magnitud comparable al de $\tau=1$.
\comprobacion cada gradiente suma~0, porque $q$ y $p$ suman~1: $-0{,}2677+0{,}0977+0{,}1699=-0{,}0001\approx 0$. Con $\tau=4$ los objetivos blandos revelan que el profesor considera el segundo token bastante más probable que el tercero ($0{,}29$ frente a $0{,}23$), información casi invisible con $\tau=1$.
\end{solucion}

\begin{solucion}{eje:5-altaT}
\datos $\partial C/\partial z_i=(q_i-p_i)/\tau$; $\abs{z_i}/\tau$ y $\abs{v_i}/\tau$ pequeños; $\sum_iz_i=\sum_iv_i=0$.
\planteamiento linealizar las exponenciales de las dos softmax~\cite[§2.1]{hinton}.
\pasos
\begin{enumerate}[label=\arabic*.]
  \item $q_i=\dfrac{e^{z_i/\tau}}{\sum_je^{z_j/\tau}}\approx\dfrac{1+z_i/\tau}{K+\sum_jz_j/\tau}=\dfrac{1+z_i/\tau}{K}$, por la media nula; análogamente $p_i\approx(1+v_i/\tau)/K$.
  \item $\dfrac{\partial C}{\partial z_i}\approx\dfrac{1}{\tau}\cdot\dfrac{z_i-v_i}{K\tau}=\dfrac{z_i-v_i}{K\tau^{2}}$. \hfill$\blacksquare$
  \item (b) La función $G(z)=\frac{1}{2K\tau^{2}}\sum_i(z_i-v_i)^{2}$ tiene gradiente $\partial G/\partial z_i=(z_i-v_i)/(K\tau^{2})$, el mismo; minimizar $C$ equivale, en ese régimen, a igualar logits por mínimos cuadrados.
  \item (c) El gradiente blando es proporcional a $1/\tau^{2}$; si se combina con la pérdida de las etiquetas (a temperatura~1) sin corregir, su peso relativo cambiaría con $\tau$. Multiplicar por $\tau^{2}$ lo mantiene aproximadamente constante.
  \item (d) Centrado: $v-\bar v=(5/3;\allowbreak\,-1/3;\allowbreak\,-4/3)$, $z-\bar z=(2/3;\allowbreak\,-1/3;\allowbreak\,-1/3)$; diferencia $(-1;\allowbreak\,0;\allowbreak\,1)$. Aproximación con $K=3$, $\tau=10$: $(-1;\allowbreak\,0;\allowbreak\,1)/300=(-0{,}00333;\allowbreak\,0;\allowbreak\,0{,}00333)$.
  \item Valor exacto con $\tau=10$ (la softmax no cambia al centrar): $p=(0{,}3907;\allowbreak\,0{,}3199;\allowbreak\,0{,}2894)$, $q=(0{,}3559;\allowbreak\,0{,}3220;\allowbreak\,0{,}3220)$, $(q-p)/10=(-0{,}00348;\allowbreak\,0{,}00022;\allowbreak\,0{,}00326)$.
\end{enumerate}
\resultado \textbf{$\partial C/\partial z_i\approx(z_i-v_i)/(K\tau^{2})$; la destilación a temperatura alta iguala logits por mínimos cuadrados, y el factor $\tau^{2}$ compensa la escala $1/\tau^{2}$.} Con $\tau=10$ la aproximación $(-0{,}00333;\allowbreak\,0;\allowbreak\,0{,}00333)$ se acerca al valor exacto $(-0{,}00348;\allowbreak\,0{,}00022;\allowbreak\,0{,}00326)$.
\comprobacion con $\tau=4$ (\cref{eje:5-destila}) la aproximación da $(-1;0;1)/48=(-0{,}0208;\allowbreak\,0;\allowbreak\,0{,}0208)$ frente a $(-0{,}0225;\allowbreak\,0{,}0032;\allowbreak\,0{,}0193)$: el error relativo disminuye al crecer $\tau$.
\end{solucion}

% ---------------------------------------------------------------------
\section{Tema 6}

\begin{solucion}{eje:6-paginacion}
\planteamiento comparar la \cref{def:paged} con la memoria virtual paginada.
\pasos
\begin{enumerate}[label=\arabic*.]
  \item Un proceso ve un espacio de direcciones contiguo formado por páginas físicas dispersas; la tabla de páginas traduce direcciones lógicas a físicas. En PagedAttention, cada secuencia ve una caché lógica contigua formada por bloques físicos dispersos, y la tabla de bloques hace la traducción.
  \item Los bloques se asignan cuando se necesitan, de modo que no hay que reservar el tamaño máximo: desaparece casi toda la fragmentación interna (solo queda parte del último bloque, \cref{prop:desperdicio}). Como todos los bloques son iguales, no hay fragmentación externa.
  \item Compartir páginas de solo lectura entre procesos equivale a que varias secuencias apunten a los mismos bloques del prompt común; la copia en escritura de los sistemas operativos corresponde a copiar el bloque parcial cuando una secuencia escribe en él (\cref{prop:prefijo}).
\end{enumerate}
\resultado \textbf{PagedAttention elimina la reserva contigua y la fragmentación, y permite compartir bloques con copia en escritura, igual que la paginación de un sistema operativo.}
\end{solucion}

\begin{solucion}{eje:6-bloques}
\datos $\ell=1000$; $b_{\mathrm{blk}}=16$; reserva contigua de $2048$.
\planteamiento \cref{prop:desperdicio}.
\pasos
\begin{enumerate}[label=\arabic*.]
  \item $1000/16=62{,}5$, luego $\lceil 62{,}5\rceil=63$ bloques.
  \item Posiciones: $63\cdot 16=1008$; desperdicio $1008-1000=8$, es decir, $8/1008=0{,}794\,\%$.
  \item Reserva contigua: $2048-1000=1048$ posiciones vacías, $1048/2048=51{,}17\,\%$.
\end{enumerate}
\resultado \textbf{63 bloques y 8 posiciones desperdiciadas ($0{,}794\,\%$), frente a 1048 ($51{,}17\,\%$) con la reserva contigua.}
\comprobacion $8\le b_{\mathrm{blk}}-1=15$, como garantiza la \cref{prop:desperdicio}.
\end{solucion}

\begin{solucion}{eje:6-latencia}
\datos $N=1{,}3\cdot 10^{10}$; fp16; $1{,}5\cdot 10^{14}$~FLOP/s efectivos en el \emph{prefill}; $\Bmem=10^{12}$~bytes/s; $n=2000$; $n_{\mathrm{out}}=500$; longitudes $100$, $400$, $250$, $50$.
\planteamiento $\TTFT\approx 2Nn/(\text{FLOP/s efectivos})$ (\cref{ej:prefill}); $\TPOT\approx sN/\Bmem$ (\cref{teo:cota-memoria}); latencia por la \cref{prop:latencia}; utilización por la \cref{prop:utilizacion}.
\pasos
\begin{enumerate}[label=\arabic*.]
  \item $\TTFT$: $2\cdot 1{,}3\cdot 10^{10}\cdot 2000=5{,}2\cdot 10^{13}$~FLOP; $5{,}2\cdot 10^{13}/1{,}5\cdot 10^{14}=0{,}3467$~s.
  \item $\TPOT$: $2{,}6\cdot 10^{10}/10^{12}=0{,}026$~s ($38{,}46$ tokens/s).
  \item Latencia: $0{,}3467+499\cdot 0{,}026=0{,}3467+12{,}97=13{,}32$~s.
  \item (b) $\sum\ell_i=800$; $B\max\ell_i=4\cdot 400=1600$; $U=800/1600=0{,}5$.
\end{enumerate}
\resultado \textbf{(a) $\TTFT\approx 0{,}347$~s, $\TPOT\approx 26$~ms y latencia total $\approx 13{,}32$~s. (b) Utilización del $50\,\%$.}
\comprobacion la decodificación domina la latencia ($12{,}97$ de $13{,}32$~s) aunque el \emph{prefill} procese cuatro veces más tokens, porque cada paso de decodificación está limitado por la lectura de los $26$~GB de pesos.
\end{solucion}

\begin{solucion}{eje:6-gqa}
\datos $N'=1{,}1338\cdot 10^{10}$; bf16; $L=40$, $\nkv=8$, $\dhd=128$; $S=1024$; $\Bmem=3{,}35\cdot 10^{12}$; 80~GB.
\planteamiento los mismos cálculos de la solución del \cref{eje:2-llama13} con $\nkv=8$.
\pasos
\begin{enumerate}[label=\arabic*.]
  \item $c_{\mathrm{kv}}=2\cdot 40\cdot 8\cdot 128\cdot 2=163\,840$~bytes; por secuencia, $1024\cdot 163\,840=1{,}678\cdot 10^{8}$~bytes $=0{,}1678$~GB, cinco veces menos.
  \item Pesos: $2\cdot 1{,}1338\cdot 10^{10}=22{,}68$~GB.
  \item $B=64$: caché $10{,}74$~GB; memoria $33{,}41$~GB; latencia $33{,}41\cdot 10^{9}/3{,}35\cdot 10^{12}=9{,}97$~ms; rendimiento $64/0{,}00997=6417$ tokens/s.
  \item $B=256$: caché $42{,}95$~GB; memoria $65{,}63$~GB; latencia $19{,}59$~ms; rendimiento $256/0{,}01959=13\,068$ tokens/s.
  \item Ambos caben en 80~GB. Con atención multicabeza, $B=64$ daba $23{,}80$~ms y $2689$ tokens/s, y $B=256$ no cabía.
\end{enumerate}
\resultado \textbf{caché de $0{,}168$~GB por secuencia; $B=64$: $33{,}41$~GB, $9{,}97$~ms, $6417$ tokens/s; $B=256$: $65{,}63$~GB, $19{,}59$~ms, $13\,068$ tokens/s.} GQA permite lotes cuatro veces mayores en la misma GPU y multiplica el rendimiento por casi cinco respecto al mejor caso que cabía con atención multicabeza.
\comprobacion la caché se divide exactamente por $40/8=5$: $0{,}8389/0{,}1678=5{,}0$.
\end{solucion}

\begin{solucion}{eje:6-prefijo}
\datos $n_s=4$; $P=1000$; $b_{\mathrm{blk}}=16$; $512$~KiB por token.
\planteamiento \cref{prop:prefijo}.
\pasos
\begin{enumerate}[label=\arabic*.]
  \item Sin compartir: cada muestra necesita $\lceil 1000/16\rceil=63$ bloques; en total $4\cdot 63=252$.
  \item Con compartición: $\lfloor 1000/16\rfloor=62$ bloques completos compartidos; el bloque parcial ($1000-62\cdot 16=8$ tokens) se copia una vez por muestra: $62+4=66$ bloques.
  \item Ahorro: $252-66=186$ bloques. Cada bloque ocupa $16\cdot 512=8192$~KiB $=8$~MiB, luego se ahorran $186\cdot 8=1488$~MiB.
\end{enumerate}
\resultado \textbf{252 bloques sin compartir frente a 66 con compartición; se ahorran 186 bloques, unos 1488~MiB ($1{,}45$~GiB).}
\comprobacion el ahorro es $(n_s-1)\lfloor P/b_{\mathrm{blk}}\rfloor=3\cdot 62=186$ bloques: cada muestra adicional reutiliza los 62 bloques completos.
\end{solucion}

\begin{solucion}{eje:6-frag}
\datos $r=\ell\bmod b_{\mathrm{blk}}$ uniforme en $\{0,\dots,b_{\mathrm{blk}}-1\}$; en (b), $\ell$ uniforme en $\{1,\dots,L_{\max}\}$.
\planteamiento expresar el desperdicio en función de $r$ o de $\ell$ y tomar esperanzas.
\pasos
\begin{enumerate}[label=\arabic*.]
  \item (a) Por la \cref{prop:desperdicio}, el desperdicio es $w=b_{\mathrm{blk}}\lceil\ell/b_{\mathrm{blk}}\rceil-\ell$, que vale $0$ si $r=0$ y $b_{\mathrm{blk}}-r$ si $r\in\{1,\dots,b_{\mathrm{blk}}-1\}$.
  \item Cuando $r$ recorre $\{0,\dots,b_{\mathrm{blk}}-1\}$, $w$ recorre $\{0\}\cup\{b_{\mathrm{blk}}-1,\dots,1\}=\{0,\dots,b_{\mathrm{blk}}-1\}$, cada valor una vez; luego $w$ es uniforme en ese conjunto.
  \item $\E[w]=\frac{1}{b_{\mathrm{blk}}}\sum_{j=0}^{b_{\mathrm{blk}}-1}j=\frac{1}{b_{\mathrm{blk}}}\cdot\frac{(b_{\mathrm{blk}}-1)b_{\mathrm{blk}}}{2}=\frac{b_{\mathrm{blk}}-1}{2}$. \hfill$\blacksquare$
  \item (b) Con reserva contigua el desperdicio es $L_{\max}-\ell$ y $\E[\ell]=(L_{\max}+1)/2$, luego $\E[L_{\max}-\ell]=(L_{\max}-1)/2$. \hfill$\blacksquare$
  \item (c) $(16-1)/2=7{,}5$ posiciones; $(2048-1)/2=1023{,}5$ posiciones.
\end{enumerate}
\resultado \textbf{desperdicio esperado de $(b_{\mathrm{blk}}-1)/2=7{,}5$ posiciones con bloques de 16, frente a $(L_{\max}-1)/2=1023{,}5$ con reserva contigua de 2048: unas 136 veces menos.}
\comprobacion $1023{,}5/7{,}5=136{,}5$. El modelo de (b) predice un uso medio del $50\,\%$ de la reserva; Kwon \emph{et al.} miden entre el $20{,}4\,\%$ y el $38{,}2\,\%$ en sistemas reales, que además reservan para el caso peor y sufren fragmentación externa~\cite{kwon}.
\end{solucion}
